%&pdftex
%%%%%%%%%%%%%%%%%%%%%%%%%%%%%%%%%%%%%%%%%%%%%%%%%%%%%%%%%%%%%%%%%%%%%%%%%%%%%%%
%             MEAN AND WEAK CONVERGENCE OF FOURIER-BESSEL SERIES,             %
%                                     por                                     %
%   Jos\'e J. Guadalupe, Mario P\'ez, Francisco J. Ruiz y Juan Luis Varona.   %
%                                                                             %
%              ULTIMA MODIFICACION: 29-ENE-93, por Mario P\'erez.             %
%                                                                             %
%      PUBLICADO EN EL JOURNAL OF MATHEMATICAL ANALYSIS AND APPLICATIONS.     %
%%%%%%%%%%%%%%%%%%%%%%%%%%%%%%%%%%%%%%%%%%%%%%%%%%%%%%%%%%%%%%%%%%%%%%%%%%%%%%%

\magnification=1200

\font\msbmnormal=msbm10
\font\msbmpeq=msbm7
\font\msbmmuypeq=msbm5
\newfam\numeros
\textfont\numeros=\msbmnormal
\scriptfont\numeros=\msbmpeq
\scriptscriptfont\numeros=\msbmmuypeq
\def\num{\fam\numeros\msbmnormal}
\def\R{{\num R}}

\font\rmdiez=cmr10 at 8pt

\def\proof{\hbox{\it Proof.\enspace}}
\def\por{\thinspace}
% \baselineskip=24pt

%%%%%%%%%%%%%%%%%%%%%%%%%%%%%%%%%%%%%%%%%%%%%%%%%%%%%%%%%%%%%%%%%%%%%
%%%%%%%%%%%%%%%%%%%%% FIRST PAGES WITH INFORMATION %%%%%%%%%%%%%%%%%%
%%%%%%%%%%%%%%%%%%%%%%%%%%%%%%%%%%%%%%%%%%%%%%%%%%%%%%%%%%%%%%%%%%%%%
%	\centerline{\bf MEAN AND WEAK CONVERGENCE OF FOURIER-BESSEL SERIES}  
%	\medskip
%	\centerline {by}
%	\medskip
%	\centerline {J. J. GUADALUPE, M. PEREZ, F. J. RUIZ and J. L. VARONA}
%	\footnote{}{AMS (1991) Subject Classification: {\bf 42C10}}
%	\footnote{}{Keywords and phrases: Convergence of Fourier series, Bessel
%	functions, weighted norm}
%	\footnote{}{inequalities, $A_p$-weights, weak and restricted weak type.}
%	\footnote{}{Research supported by DGICYT under grant PB89-0181-C02-02.}
%	\bigskip
%
%	Jos\'e J. Guadalupe, Mario P\'erez, Francisco
%	J. Ruiz. Departamento de Matem\'aticas. Facultad de Ciencias.
%	Universidad de Zaragoza, 50009 Zaragoza, SPAIN.
%
%	\bigskip
%
%	Juan L. Varona. Departamento de Matem\'atica
%	Aplicada. Colegio Universitario de la Rioja, 26001 Logro\~no, SPAIN.
%
%	\vfill
%	\eject
%
%	Proposed running head: Fourier-Bessel series.
%
%	\bigskip
%
%	Mailing address:
%
%	\centerline{Francisco J. Ruiz}
%	\centerline{Departamento de Matem\'aticas.}
%	\centerline{Facultad de Ciencias.}
%	\centerline{Universidad de Zaragoza.}
%	\centerline{50009 Zaragoza.}
%	\centerline{Spain.}
%
%	e-mail: JVARONA@CC.UNIZAR.ES
%
%	\vfill
%	\eject
%%%%%%%%%%%%%%%%%%%%%%%%%%%%%%%%%%%%%%%%%%%%%%%%%%%%%%%%%%%%%%%%%%%%%

\centerline{\bf MEAN AND WEAK CONVERGENCE OF FOURIER-BESSEL SERIES}  
\medskip
\centerline {by}
\medskip
\centerline {J. J. GUADALUPE, M. PEREZ, F. J. RUIZ and J. L. VARONA}
\footnote{}{AMS (1991) Subject Classification: {\bf 42C10}}
\footnote{}{Keywords and phrases: Convergence of Fourier series, Bessel
functions, weighted norm}
\footnote{}{inequalities, $A_p$-weights, weak and restricted weak type.}
\footnote{}{Research supported by DGICYT under grant PB89-0181-C02-02.}
\footnote{}{THIS PAPER HAS BEEN PUBLISHED IN: 
{\it J.\ Math.\ Anal.\ Appl.}\ {\bf 173} (1993), 370--389.}
\bigskip

\midinsert
\narrower\narrower
\noindent ABSTRACT: {\sl We study the uniform boundedness on some weighted $L^p$
spaces of the partial sum operators associated to Fourier-Bessel series,
obtaining necessary and sufficient conditions for this boundedness in terms
of the weights. On the other hand, we study the weak and restricted weak
type of the partial sum operators in the end points of the interval of
strong  boundedness.}
\endinsert

\bigskip
\noindent {\bf \S 0. Introduction.}\par\nobreak
\medskip
Let $J_{\alpha}$ be the Bessel function of order $\alpha >-1$, and  
$j_n^{\alpha} (x) =  2^{1/2} |J_{\alpha+1}(\alpha_n)|^{-1}\por J_{\alpha}
(\alpha_nx)$ (where $\{\alpha_n\}$ is the increasing sequence of the zeroes of
$J_{\alpha}$)  the Bessel system of order $\alpha$, which is orthonormal
and complete in $L^2 ((0,1);xdx)$, and therefore
the Fourier series of a function $f\in L^2 ((0,1);xdx)$ with respect to this
system converges to $f$ in the $L^2$-norm.

The next step is to ask for which $p\in(1,\infty),\  p\ne 2$ the above
convergence is true. The problem can be posed, by the Banach-Steinhauss theorem,
in terms of the uniform boundedness on $L^p ((0,1);xdx)$ of the partial sum
operators 
$$
S_n^\alpha f(x) = \sum_{k=1}^n c_k j_k^\alpha (x)\ , \  \  c_k = c_k (f) =
\int_0^1 j_k^\alpha (y) f(y) y dy.
$$
The problem was solved by Wing {\bf [13]} (in the case $\alpha \geq -1/2$) and
Benedek and Panzone {\bf [1]} (general case). Moreover, these authors obtained
some weighted norm estimates of the form
$$
\Vert S_n^\alpha f(x)U(x)\Vert_p \leq C \Vert f(x)V(x)\Vert_p 
\eqno (0.1)
$$
where $C$ is a constant independent of $n$, $U(x)=V(x)$ are functions of radial
type $x^{a}$ and by $\Vert.\Vert_p$ we denote the $L^p ((0,1);xdx)$-norm.

The first objective of this paper is to extend inequality (0.1) to a wider
class of weights $U(x)$ and $V(x)$, not necessarily the same. We obtain
sufficient conditions in terms of $A_p$-conditions (propositions 1 and 2) and we
solve completely the problem (theorem 1) when the weights are of the form 
$$
U(x)=x^{a}(1-x)^{b}\prod_{k=1}^m |x-x_k|^{b_k} ;\  0<x_1<...<x_m<1;\
a,b,b_k\in\R. \eqno (0.2)
$$
The inequality (0.1) is closely related to the problem of the boundedness of the
Hilbert transform with two weights. In this sense, the use of pairs of weights
in the Muckenhoupt $A_p$ classes (see {\bf [9]}, {\bf [10]}) will be one of the
main ingredients in the proofs. We obtain also checkable necessary conditions of
integrability on $U,V$ from the inequality (0.1) (theorem 2) in the spirit of
M\'at\'e, Nevai and Totik's conditions (see {\bf [8]}) for systems of orthogonal
polynomials on $[-1,1]$.

On the other hand, we shall study a problem related to the end points of the
interval of mean convergence: by interpolation, the range of $p^,s$ such that
the uniform boundedness (0.1) holds is always an interval. When $U(x)=V(x)=1$
(for example), this interval is open, say $(p_0,p_1)$, and then the natural
question is if  $S_n^\alpha$ are ``uniformly" of weak type $(p_0,p_0)$ or/and 
$(p_1,p_1)$. That is, if we denote
$$
\Vert g\Vert_{L^p_* ((0,1);U(x)^pxdx)}= \{\displaystyle\sup_{\lambda >0}
\lambda^p\int_{\{x:|g(x)|>\lambda\}}U(x)^pxdx\}^{1/p}
$$
then the question is if
$$
\Vert S_n^\alpha f\Vert_{L^p_* ((0,1);U(x)^pxdx)}\leq C\Vert f(x)V(x) \Vert_p
\eqno (0.3)
$$
with  a constant $C$ independent of $n$, and $p=p_0$ or $p_1$.

This ``end point problem" has been studied by different authors for another
operators (see {\bf [2]}, {\bf [3]}, {\bf [6]}, {\bf [11]}), but in the case of
Fourier-Bessel series it seems to be new, even in the unweighted case. In
theorem 3, we give necessary conditions on $U,V$ for (0.3) to be true. In the
particular case $\alpha\geq -1/2$ and $U(x)=V(x)$ of the form (0.2), we show
that (0.3) is false for both end points (theorem 4). Finally, when  $\alpha\geq
-1/2$ and $U(x)=V(x)=1$ we show (theorem 5) that $S_n^\alpha$ are ``uniformly" of
restricted weak type at the end points, i.e., (0.3) is true, but only for
functions $f=\chi_E$ where $E$ is a measurable set contained in $(0,1)$.

The organization of the paper is as follows: in section 1, we exhibit the main
results with the notations which are needed to make them understandable. In
section 2, we collect several technical lemmas concerning $A_p$ weights and
the boundedness of some operators associated to the problem which we shall use
for the proofs of the main results which will be given in section 3.

\bigskip
\noindent {\bf \S 1. Main results.}\par\nobreak
\medskip
Let $1 < p < \infty$ and $-\infty\leq a < b \leq \infty$. $A_p(a,b)$ will
stand for the Muckenhoupt classes (see {\bf [4]}, {\bf [9]}) consisting of 
those pairs of nonnegative functions $(u,v)$ on $(a,b)$ such that 
$$
 (|I|^{-1} \int_I u ) (|I|^{-1} \int_I v^{-p'/p})^{p/p'}
\leq C \quad {\rm  for\  every\  interval \ } I\subseteq (a,b),
$$
where $p+p'=pp'$.

We say that $(u,v)\in A_p^\delta(a,b)$ (where $\delta >1$) if
$(u^\delta,v^\delta)\in A_p(a,b)$. It is clear from Holder's inequality
that $A_p^\delta(a,b)\subseteq A_p(a,b)$.

We say that a sequence $\{(u_n,v_n)\}_{n=1}^\infty \in A_p^\delta(a,b)$
``uniformly" if $(u_n,v_n)\in A_p^\delta(a,b)$ , $\forall n$ and the constant
appearing in the definition is independent of $n$.

In the inequality \ $\Vert S_n^\alpha f(x)U(x)\Vert_p \leq C \Vert
f(x)V(x)\Vert_p$\  we shall always mean that $C$ is a constant independent of $f$
and $n$ (and the same for the weak $L^p_* $-norm).

In our two first results, we distinguish the cases $\alpha\geq -1/2$ and
$-1<\alpha <-1/2$. The reason for that is the different assymptotic behavior of
the Bessel functions in each case.

\proclaim {Proposition 1}. (Case $\alpha\geq -1/2$). Let $1 < p < \infty$. Let
$U$, $V$ be weights on $(0,1)$. If $(U(x)^px^{1-p/2},V(x)^px^{1-p/2})\in
A_p^\delta(0,1)$ for some $\delta > 1$ ($\delta = 1$ if $U=V$), then
$$
   \Vert S_n^\alpha f(x)U(x)\Vert_p\leq C \Vert f(x)V(x)\Vert_p.
$$

\proclaim {Proposition 2}. (Case $-1<\alpha <-1/2$). Let  $1 < p < \infty$ and
$M_n=(\alpha_n+\alpha_{n+1})/2$. Let $U$, $V$ be weights on $(0,1)$. If there
exists $\delta > 1$ ($\delta = 1$ if $U=V$) such that
$$
   ((M_n^{-1}+x)^{-p(\alpha+1/2)}U(x)^px^{\alpha p+1},
   (M_n^{-1}+x)^{-p(\alpha+1/2)}V(x)^px^{\alpha p+1})\in A_p^\delta(0,1),
   \eqno(1.1)
$$
$$
   ((M_n^{-1}+x)^{p(\alpha+1/2)}U(x)^px^{-\alpha p-p+1},
   (M_n^{-1}+x)^{p(\alpha+1/2)}V(x)^px^{-\alpha p-p+1})\in A_p^\delta(0,1)
   \eqno(1.2)
$$
``uniformly", then \ $\Vert S_n^\alpha f(x)U(x)\Vert_p \leq C \Vert
f(x)V(x)\Vert_p$.

\noindent {\bf Note.} When $U=V$ it suffices $\delta = 1$ because from reverse
Holder's inequality it follows that \ $(u,u)\in A_p \Rightarrow (u,u)\in
A_p^\delta$ for some $\delta >1$ (see {\bf [4]}).

Applying both propositions to particular weights of the form (0.2), we have

\proclaim {Theorem 1}. Let $\alpha > -1$, $1 < p < \infty$ and the weights
$$
U(x)=x^{a}(1-x)^{b}\prod_{k=1}^m |x-x_k|^{b_k} ,\quad
V(x)=x^{A}(1-x)^{B}\prod_{k=1}^m |x-x_k|^{B_k},
$$
where  $0<x_1<...<x_m<1$ and  $a,A,b,B,b_k,B_k\in\R$. If the conditions:
$$
   \eqalign{
   &B \leq b, \qquad p B < p-1, \qquad -1 < p b; \cr
   &B_k \leq b_k, \qquad p B_k < p-1, \qquad -1 < p b_k \qquad (1 \leq k \leq m);
   \cr} \eqno(1.3)
$$
$$
   A \leq a, \qquad
   \vert {1 \over p}-{1 \over 2}+{a \over 4}+{A \over 4} \vert < {a-A \over 4} +
   \min \{{1 \over 4},{\alpha +1 \over 2} \} 
   \eqno(1.4)
$$
are fullfilled, then \ $\Vert S_n^\alpha f(x)U(x)\Vert_p \leq C \Vert
f(x)V(x)\Vert_p$.

The following general result about necessary conditions will imply, in
particular, that (1.3), (1.4) are best possible.

\proclaim {Theorem 2}. If the inequality $\Vert S_n^\alpha f(x)U(x)\Vert_p
\leq  C \Vert f(x)V(x)\Vert_p$ holds, then $U$, $V$ must satisfy the conditions
$$
   \int_0^1 U(x)^px^{\alpha p+1}dx<+\infty, \qquad
   \int_0^1 U(x)^px^{1-p/2}dx<+\infty
   \eqno(1.5)
$$
$$
   \int_0^1 V(x)^{-p'}x^{\alpha p'+1}dx<+\infty, \qquad
   \int_0^1 V(x)^{-p'}x^{1-p'/2}dx<+\infty
   \eqno(1.6)
$$
$$
   U(x) \leq C\por V(x)\  a.e.
   \eqno(1.7)
$$

\proclaim {Corollary 1}. With the same notation as in theorem 1, the
conditions (1.3) and (1.4) are also necessary for the uniform boundedness
$ \Vert S_n^\alpha f(x)U(x)\Vert_p \leq C \Vert f(x)V(x)\Vert_p $.

Regarding weak boundedness, we  begin with a version of
theorem 2 in this case:

\proclaim {Theorem 3}. Let $\alpha > -1$ and $1<p<\infty$. If the inequality
$$
\Vert S_n^\alpha f\Vert_{L^p_* ((0,1);U(x)^pxdx)}\leq C\Vert f(x)V(x) \Vert_p
$$
holds, then $U$, $V$ must satisfy conditions (1.6), (1.7) and
$$
\displaystyle\sup_{\lambda >0}
\lambda^p\int_{\{x:x^\alpha>\lambda\}}U(x)^pxdx < +\infty ,\quad
\displaystyle\sup_{\lambda >0}
\lambda^p\int_{\{x:x^{-1/2}>\lambda\}}U(x)^pxdx < +\infty.\eqno (1.5)'
$$

From this theorem, straightforward calculations show that the weak boundedness
(0.3) fails in the lower end point of the interval of strong boundedness
given by (1.3) and (1.4). However, theorem 3 does not work in the upper end
point. In this case, we have obtained that the weak boundedness is also false
(in the particular case $\alpha \geq -1/2$ and $U(x)=V(x)=$ $
x^{a}(1-x)^{b}\prod_{k=1}^m |x-x_k|^{b_k}$), but the proof is more intricate,
being necessary a more detailed analysis of the  partial sum operators.
Something very similar occurs when one considers Jacobi-Fourier series (see {\bf
[2]}, {\bf [6]}). The result can be stated as follows:

\proclaim {Theorem 4}. Let $\alpha \geq -1/2$ , $1<p<\infty$ and 
$U(x)= x^{a}(1-x)^{b}\prod_{k=1}^m|x-x_k|^{b_k}$. Then, the
following are equivalent:
\item {(i)} $\Vert S_n^\alpha f\Vert_{L^p_* ((0,1);U(x)^pxdx)}\leq C\Vert
f(x)U(x)\Vert_p$
\item {(ii)} $ \vert {1 \over p}+{a-1 \over 2}\vert < {1 \over 4};\ -1<pb<p-1;\
-1<pb_k<p-1\ (1\leq k \leq m)$.

\medskip
Finally, following the method of {\bf [2, {\rm theorem} 1]}, we establish this
(p,p)-restricted weak type:

\proclaim {Theorem 5}. Let $\alpha \geq -1/2$ and $p=4$ or $p=4/3$.Then
$$
   \Vert S_n^\alpha (\chi_E)\Vert_{L^p_* ((0,1);xdx)}\leq C \Vert
   \chi_E\Vert_p
$$
for every measurable set $E\subseteq (0,1)$.

\medskip
\noindent {\bf \S 2. Technical lemmas.}\par\nobreak
\medskip
We denote by $K_n^\alpha(x,y) = \sum_{k=1}^n j_k^\alpha (x)j_k^\alpha(y)$ the
kernel of the operator $S_n^\alpha$, that is, 
$$
S_n^\alpha f(x) = \int_0^1 K_n^\alpha(x,y) f(y)ydy.
$$
It is known the following decomposition of the kernel (see {\bf [1]}, {\bf
[13]}): $$
\eqalign{K_n^\alpha(x,y) &= J_\alpha(M_nx)J_{\alpha+1}(M_ny) {M_n \over 2(y-x)}+
J_\alpha(M_ny)J_{\alpha+1}(M_nx) {M_n \over 2(x-y)}+ \cr
&+ J_\alpha(M_nx)J_{\alpha+1}(M_ny) {M_n \over 2(y+x)}+
J_\alpha(M_ny)J_{\alpha+1}(M_nx) {M_n \over 2(x+y)}+ \cr
&+ O(1) (xy)^\alpha + O(1) { (xy)^{-1/2} \over 2-x-y} =
\sum_{i=1}^6K_{n,i}^\alpha(x,y).}
$$
So, we can decompose
$$
S_n^\alpha f(x) = \sum_{i=1}^6 S_{n,i}^\alpha f(x) =
\sum_{i=1}^6 \int_0^1 K_{n,i}^\alpha(x,y)f(y)ydy.
$$
We shall use the following notations
$$
Hf(x)=p.v.\int_0^1 {f(y) \over x-y}dy;\quad Jf(x)=\int_0^1 {f(y) \over 2-x-y}dy
;\quad J'f(x)=\int_0^1 {f(y) \over x+y}dy. 
$$
 The main tools for estimating the operators $S_{n,i}^\alpha f$ are:
$$
   J_\alpha(z)=z^\alpha2^{-\alpha}\Gamma(\alpha+1)^{-1}+O(z^{\alpha+2}),\quad
   z\longrightarrow0
   \eqno(2.1)
$$
$$
   J_\alpha(z)= 2^{1/2} (\pi z)^{-1/2}
   \cos (z-{ \alpha\pi \over 2}-{ \pi\over 4}) + O(z^{-3/2})
   \quad z\longrightarrow\infty
   \eqno(2.2)
$$
$$
   \int_0^1 J_\alpha(\alpha_nx)^2x\por dx \approx {1 \over \pi\alpha_n},
   \quad n\longrightarrow \infty
   \eqno(2.3)
$$
for $\alpha > -1$ (see {\bf [12]}) and the following results:

\proclaim {Lemma 1 }. Let $1<p<\infty$. If $(u,v)\in
A_p^\delta(0,1)$ for some $\delta >1$, then
$$
\Vert Tf(x)\Vert_{L^p ((0,1);u(x)dx)}\leq 
C \Vert f(x)\Vert_{L^p ((0,1);v(x)dx)}.
$$
where $T=H,\ J$ or $J'$.

\proof The result for $H$ is due to Neugebauer {\bf[10]}. In fact, Neugebauer
showed that under the hypothesis there exists a $w\in A_p$ (i.e. $(w,w)\in A_p$)
and constants $C_1, C_2$ such that $C_1u(x)\leq w(x)\leq C_2v(x)$ a.e., and then
we can apply that $H$ is bounded on $L^p ((0,1);w(x)dx)$ if and only if $w\in
A_p$ (see {\bf[7]}). The last fact is also true for $Mf$, being $M$ the
Hardy-Littlewood maximal operator (see {\bf[9]}).

Let now  $T_1f(x)={1 \over x}\int_0^x f(y)dy$ the Hardy function and
$T_2f(x)=\int_x^1{f(y) \over y}dy$ its adjoint. Then, it is easy to see that
$$
|J'f(x)|\leq T_1(|f|)(x) + T_2(|f|)(x);\quad|T_1f(x)|\leq 2J'(|f|)(x) ;
\quad|T_2f(x)|\leq 2J'(|f|)(x).
$$
So, the weighted boundedness for $J'$ is equivalent to the one for $T_1$ and
$T_2$. The operator  $T_1$ is controlled by $M$; and for its adjoint $T_2$ we use
the duality property for $A_p$-classes: $(u,v)\in A_p$ if and only if
$(v^{-p'/p},u^{-p'/p})\in A_{p'}$. So, $J'$ is controlled by operators which
are bounded under the hypothesis.

Finally, we can deal with the operator $J$ in the same way as  $J'$ by making the
change of variables $X=1-x, Y=1-y$ which does not change $A_p$ conditions.
\medskip
\noindent {\bf Note.} The constants $``C"$ appearing in the weighted norm
inequalities  depend only on the constants appearing in the definition of the
$A_p$ classes. Thus, for sequences of weights belonging to
$A_p^\delta$ ``uniformly", the corresponding weighted norm inequalities are
also uniform.

For the proof of theorem 1, we  need to analyze when certain sequences of
weights of radial type belong to $A_p$ classes ``uniformly". For that, we give
the 

\proclaim {Lemma 2}. Let $r,s,R,S\in \R$ , $1<p<\infty$ and let
$\{x_n\}_{n=1}^\infty$ be a sequence of positive numbers such that $\lim
\por x_n=0$. Then
$$
(x^r(x+x_n)^s , x^R(x+x_n)^S )\in A_p^\delta (0,1)\quad {\rm for \  some} \ \ 
\delta >1 \Longleftrightarrow
$$
$$
-1<r;\quad -1<r+s;\quad R<p-1;\quad R+S<p-1;\quad R\leq r;\quad R+S\leq r+s
\eqno (2.4)
$$

\proof Straightforward calculations show that for $0\leq a < b \leq 1$, $0\leq c
\leq 1$, $\gamma >-1$, $\gamma +\mu >-1$, there exists a constant $K$,
depending only on $\gamma$, $\mu$,  such that 
$$
\int_a^b x^\gamma (x+c)^\mu dx \leq \cases {Kb^{\gamma+\mu}(b-a), &if $c\leq
b$;\cr Kb^\gamma c^\mu (b-a), &if $b\leq c$. \cr}
$$
and then (2.4) implies
$$
(\int_a^b x^r (x+c)^s dx)(\int_a^b (x^R (x+c)^S)^{-p'/p} dx)^{p/p'}\leq K
(b-a)^p,
$$
that is, $(x^r (x+c)^s,x^R (x+c)^S)\in A_p (0,1)$ ``uniformly" on $c$. Now,
these pairs of weights belong to $A_p^\delta (0,1)$, for some $\delta>1$ because
the inequalities $-1<r$; $-1<r+s$; $ R<p-1$; $ R+S<p-1$ are strict.

Reciprocally, the $A_p$ condition implies that  $x^r (x+x_n)^s$ and 
$(x^R(x+x_n)^S)^{-1/p-1}$ are integrables on $(0,1)$ and  then $-1<r$, $R<p-1$.
Taking limits on $n$ we get the integrability of $x^{r+s}$ and
$(x^{R+S})^{-1/p-1}$ (even more, the pair  $(x^{r+s},x^{R+S})\in A_p
(0,1)$).This gives $-1<r+s$ and $R+S<p-1$. Finally, by using that if $(u,v) \in
A_p(0,1)$ then $u\leq v\ a.e.$, we have  $R\leq r$,\ $R+S\leq r+s$.

On the other hand, we  also need to deal with products of weights in $A_p$
classes. In this sense it is not difficult to prove

\proclaim {Lemma 3}. Let $\{u_n(x)\}_{n=0}^\infty$, $\{v_n(x)\}_{n=0}^\infty$,
$\{U_n(x)\}_{n=0}^\infty$, $\{V_n(x)\}_{n=0}^\infty$ be sequences of weights on
a finite interval $(a,b)$. Let $c\in (a,b)$ and $\epsilon >0$. Assume that there
exist positive constants $C_i ,(i=1,2)$ such that $C_1\leq U_n(x)\leq C_2$,
$C_1\leq V_n(x)\leq C_2$ on $(a,c+\epsilon)$ and $C_1\leq u_n(x)\leq C_2$,
$C_1\leq v_n(x)\leq C_2$ on $(c-\epsilon ,b)$, $\forall n$. If $\{(u_n,v_n)\}\in
A_p(a,c)$ ``uniformly" and $\{(U_n,V_n)\}\in A_p(c,b)$ ``uniformly" then
$\{(u_nU_n,v_nV_n)\}\in A_p(a,b)$ ``uniformly".

In particular, the iterated use of lemma 3 and the fact that \ $(|x-c|^\alpha ,
|x-c|^\beta )\in A_p^\delta (a,b)$ for some $\delta >1$ (where $a\leq c \leq
b$) if and only if $-1<\alpha$, $\beta < p-1$ and $\beta \leq \alpha$ gives

\proclaim {Lemma 4}. Let $\alpha > -1$ and the weights
$$
U(x)=x^{a}(1-x)^{b}\prod_{k=1}^m |x-x_k|^{b_k} ,\quad
V(x)=x^{A}(1-x)^{B}\prod_{k=1}^m |x-x_k|^{B_k}
$$
where  $0<x_1<...<x_m<1$ and  $a,A,b,B,b_k,B_k\in\R$. Then, $(U,V)\in A_p^\delta
(0,1)$ if and only if
$$
   \eqalign{
   &A \leq a, \qquad -1 < a, \qquad A < p-1; \cr
   &B \leq b, \qquad -1 < b, \qquad B < p-1; \cr
   &B_k \leq b_k, \qquad -1 < b_k, \qquad B_k < p-1 \qquad (1 \leq k \leq m). \cr
   }
$$

The following lemma will be crucial in the study of necessary conditions for
the mean and weak convergence of Fourier Bessel series. It is parallel to the
one given by M\'at\'e, Nevai and Totik {\bf [8, {\rm theorem} 2]} in the context
of orthogonal polynomials.

\proclaim {Lemma 5}. Let $\alpha>-1$. Let $h(x)$ be a Lebesgue measurable
nonnegative function on $[-1,1]$, $\{r_n\}$ a sequence such that
$\displaystyle\lim_{n \to \infty} r_n = +\infty$ and $1\leq p<\infty$. Then 
$$
\displaystyle\lim_{n \to \infty}\int_0^1 |r_n^{1/2}J_{\alpha}(r_nx)|^ph(x)dx
\geq M \por \int_0^1 h(x)\por x^{-p/2}\por dx 
$$
where $M$ is a positive constant independent of $h(x)$ and $\{r_n\}$. In
particular, (with $r_n=\alpha_n$)
$$
\displaystyle\lim_{n \to \infty}\int_0^1 |j_n^{\alpha}(x)|^ph(x)dx \geq M \por
\int_0^1 h(x)\por x^{-p/2}\por dx. \eqno (2.5)
$$

\proof It can be seen in {\bf [5]}. There, the proof is  given only in the case
$r_n=\alpha_n$ but the arguments are the same for any sequence $\{r_n\}$  such
that $\displaystyle\lim_{n \to \infty} r_n = +\infty$.

\bigskip
\noindent {\bf \S 3. Proofs of the main results.}\par\nobreak
\medskip
{\hbox{\bf Proof of proposition 1.\enspace}} We  use the estimate
$$
|J_\alpha (z)| \leq C \por z^{-1/2},\quad z\in (0,\infty),\quad \alpha \geq -1/2
\eqno (3.1)
$$
which can be deduced from (2.1) and (2.2). We are going to prove that all the
operators $S_{n,i}^\alpha$ ($i=1$ to $6$) are uniformly bounded.
$$
\int_0^1 |S_{n,1}^\alpha f(x)U(x)|^pxdx =
$$
$$
=\int_0^1 \vert \int_0^1 { f(y)y^{1/2}(yM_n)^{1/2}J_{\alpha +1}(yM_n) \over
2(y-x)}dy\vert ^p \vert (xM_n)^{1/2}J_{\alpha}(xM_n)\vert ^p U(x)^px^{1-p/2}dx
\leq
$$
$$
\leq C\por \int_0^1\vert \int_0^1{ f(y)y^{1/2}(yM_n)^{1/2}J_{\alpha +1}(yM_n)
\over 2(y-x)}dy \vert ^p U(x)^px^{1-p/2}dx \leq
$$
$$
\leq C\por \int_0^1 |f(x)x^{1/2}(xM_n)^{1/2}J_{\alpha +1}(xM_n)|^p
V(x)^px^{1-p/2}dx \leq C \por \int_0^1 |f(x)V(x)|^p xdx,
$$
where we have used (3.1), the hypothesis on $U,V$ and lemma 1 for the Hilbert
transform.

The boundedness of $S_{n,2}^\alpha$ is completely similar and the operators
$S_{n,3}^\alpha$, $S_{n,4}^\alpha$ can be treated as the previous operators but
using the operator $J'$ instead of the Hilbert transform.

For the boundedness of $S_{n,5}^\alpha$ we only need Holder's inequality and
the fact that $u$ and $v^{-1/(p-1)}$ are integrable when $(u,v)\in A_p(0,1)$.
Finally, it is immediate that
$$
\Vert S_{n,6}^\alpha f(x)U(x)\Vert_p \leq C \Vert f(x)V(x)\Vert_p
$$ 
if and only if 
$$
\Vert Jf(x)\Vert_{L^p ((0,1);U(x)^px^{1-p/2}dx)}\leq 
C \Vert f(x)\Vert_{L^p ((0,1);V(x)^px^{1-p/2}dx)}
$$
and then we can apply lemma 1 for $J$.
\medskip
{\hbox{\bf Proof of proposition 2.\enspace}} Here, we shall use that there exists
a constant $C$ such that $\forall z \in (0,\infty)$
$$
|z^{1/2}J_\alpha (z)| \leq C \por (1+z^{\alpha+1/2}), \quad
|z^{1/2}J_{\alpha +1} (z)| \leq C \por (1+z^{\alpha+1/2})^{-1}
\eqno (3.2)
$$
which follows from (2.1) and (2.2). For $\alpha < -1/2$, we have 
$$
1+(xM_n)^{\alpha+1/2} \approx (xM_n)^{\alpha+1/2}\por (1+xM_n)^{-\alpha-1/2}
$$
and so, condition (1.1) is equivalent to
$$
((1+(xM_n)^{\alpha+1/2})^p U(x)^px^{1-p/2}, 
(1+(xM_n)^{\alpha+1/2})^p V(x)^px^{1-p/2})\in A_p^\delta (0,1)
$$
``uniformly". Now, applying (3.2) and lemma 1 we get
$$
\int_0^1 |S_{n,1}^\alpha f(x)U(x)|^pxdx =
$$
$$
=\int_0^1 \vert \int_0^1 { f(y)y^{1/2}(yM_n)^{1/2}J_{\alpha +1}(yM_n) \over
2(y-x)}dy\vert ^p \vert (xM_n)^{1/2}J_{\alpha}(xM_n)\vert ^p U(x)^px^{1-p/2}dx
\leq
$$
$$
\leq C \por \int_0^1 \vert \int_0^1 { f(y)y^{1/2}(yM_n)^{1/2}J_{\alpha +1}(yM_n)
\over 2(y-x)}dy\vert ^p (1+(xM_n)^{\alpha+1/2})^p U(x)^px^{1-p/2}dx \leq
$$
$$
\leq C\por \int_0^1 |f(x)x^{1/2}(xM_n)^{1/2}J_{\alpha +1}(xM_n)|^p
(1+(xM_n)^{\alpha+1/2})^p V(x)^px^{1-p/2}dx \leq 
$$
$$
\leq C \por \int_0^1 |f(x)V(x)|^pxdx.
$$
The boundedness of $S_{n,2}^\alpha$ is similar since (1.2) is equivalent to
$$
((1+(xM_n)^{\alpha+1/2})^{-p} U(x)^px^{1-p/2}, 
(1+(xM_n)^{\alpha+1/2})^{-p} V(x)^px^{1-p/2})\in A_p^\delta (0,1).
$$
The boundedness of  $S_{n,3}^\alpha$ and $S_{n,4}^\alpha$ is analogous, 
taking the operator $J'$ instead of $H$. From conditions (1.1) and (1.2)
(taking $n=1$) it follows that
$$
\int_0^1 U(x)^px^{\alpha p+1}dx < + \infty\quad, \quad
\int_0^1 V(x)^{-p'}x^{\alpha p'+1}dx < + \infty
$$
and this fact together with Holder's inequality leads to the boundedness of
$S_{n,5}^\alpha$.

Finally, taking limits on $n$ in condition (1.1) we get 
$(U(x)^p x^{1-p/2}, V(x)^p x^{1-p/2})\in A_p^\delta (0,1)$ and then (
see the proof of proposition 1), $S_{n,6}^\alpha$ is bounded.
\medskip
{\hbox{\bf Proof of theorem 1.\enspace}} For $\alpha \geq -1/2$ , the result
follows from proposition 1 and lemma 4. Let $-1< \alpha < -1/2$; we must check
conditions (1.1) and (1.2). If we put $U(x) = x^a U_1(x)$ , $V(x) = x^A V_1(x)$,
then $U_1$ and $V_1$ are  constant-like near 0. Thus, lemma 3 tells us that
we just need to verify the conditions
$$
(U_1(x)^p,V_1(x)^p)\in A_p^\delta (0,1),
$$
$$
((x+{1\over M_n})^{ p(\alpha+1/2)} x^{-\alpha p-p+1+ap},
(x+{1\over M_n})^{ p(\alpha+1/2)} x^{-\alpha p-p+1+Ap})\in A_p^\delta (0,1),
$$
$$
((x+{1\over M_n})^{-p(\alpha+1/2)} x^{\alpha p+1+ap},
(x+{1\over M_n})^{-p(\alpha+1/2)} x^{\alpha p+1+Ap})\in A_p^\delta (0,1).
$$
With the help of lemma 4 we obtain the first condition  and
the two last ones are an easy consequence of lemma 2.
\medskip
{\hbox{\bf Proof of theorem 2.\enspace}} Let
$T_nf=S_nf-S_{n-1}f=c_n(f)j_n^\alpha$. From the hypothesis, we have 
$$
|c_n(f)|\por \Vert j_n^\alpha(x)U(x) \Vert_p \leq C \Vert f(x)V(x) \Vert_p
$$
with constant independent of $n$. Moreover, by applying duality to the
operator 

\noindent $f \in L^p((0,1);V(x)^pxdx)\ \longrightarrow c_n(f)$ \  we obtain
$$
\Vert j_n^\alpha(x)U(x) \Vert_p \por \Vert j_n^\alpha(x)V(x)^{-1}
\Vert_{p'} \leq C \eqno (3.3)
$$
``uniformly" on $n$. Now, by using (2.5) in this inequality we get the two
second conditions in (1.5) and (1.6).

On the other hand, taking $n=1$ in (3.3) and using that
$J_\alpha(\alpha_1x)\geq C x^\alpha$ when $x\in (0,1/2)$ (which follows from
(2.1)), we find the necessary conditions
$$
\int_0^{1/2} U(x)^p x^{\alpha p+1}dx<+\infty, \ \ 
\int_0^{1/2} V(x)^{-p'}x^{\alpha p'+1}dx<+\infty
$$
which combined with the two already obtained conditions gives (1.5) and (1.6).

Finally, we take the function $f(x) = \min
\{U(x),V(x)^{-1},U(x)V(x)^{-1}\}j_0(x)$. By using Holder's inequality and (3.3)
it is easy to see that $f\in  L^p((0,1);V(x)^pxdx) \cap L^2((0,1);xdx)$. For
every measurable set $E\subset (0,1)$, let $g(x)=f(x) \chi_E (x)$. Since
$S_ng \longrightarrow g$ in $L^2(xdx)$, there exists a subsequence such
that $S_{n_j}g \longrightarrow g$ a.e.. Then, by using Fatou's lemma and the
uniform boundedness of $\{S_{n_j}\}$ from $L^p((0,1);V(x)^pxdx)$ into
$L^p((0,1);U(x)^pxdx)$ we obtain
$$
\int_E |f(x)|^pU(x)^p x dx = \int_0^1 |g(x)|^pU(x)^p x dx \leq
\liminf \int_0^1 |S_{n_j}g(x)|^pU(x)^p x dx \leq
$$
$$
\leq C \int_0^1 |g(x)|^pV(x)^p x dx = C \int_E |f(x)|^pV(x)^p x dx
$$
$\forall E\subset (0,1)$, which implies (1.7) easily.
\medskip
{\hbox{\bf Proof of theorem 3.\enspace}} As in the proof of theorem 2, by
applying the uniform boundedness of $T_n$ we get
$$
\Vert j_n^\alpha(x)\Vert_{L^p_* ((0,1);U(x)^pxdx)} 
\Vert j_n^\alpha(x)V(x)^{-1}\Vert_{p'} \leq C \eqno (3.4)
$$
with $C$ independent of $n$.

We shall use the following ``Kolmogorov's inequality":
$$
\Vert g\Vert_{L^p_* (h(x)dx)} \leq
\displaystyle\sup_{E} \Vert g\chi_E\Vert_{L^r (h(x)dx)}
\Vert \chi_E\Vert_{L^s (h(x)dx)}^{-1} \leq ({p \over p-r})^{1/r}
\Vert g\Vert_{L^p_* (h(x)dx)}
$$
where $h(x)dx$ is a Borel measure, $0<r<p<\infty$, $1/s=1/r-1/p$ and the
supremum is taken over all measurable sets $E$ of finite, positive Lebesgue
measure  (see {\bf [4, {\rm lemma} V.2.8, {\rm p.} 485]}).

This result and (2.5) imply the existence of a positive constant $C$
independent of $h$ and $n$ such that
$$
\displaystyle\liminf_{n\rightarrow \infty} \Vert j_n^\alpha(x)
\Vert_{L^p_* (h(x)dx)} \geq C \Vert x^{-1/2}\Vert_{L^p_* (h(x)dx)}.
$$
By applying this condition and (2.5) again in (3.4) we obtain the second
conditions in (1.5)' and (1.6). 

In order to obtain the two other conditions in (1.5)' and (1.6), we observe
firstly that it suffices to prove
$$
\Vert x^\alpha \chi_{(0,1/2)}\Vert_{L^p_* ((0,1);U(x)^pxdx)} < \infty,\quad
\int_0^{1/2} V(x)^{-p'}x^{\alpha p'+1}dx<+\infty \eqno (3.5)
$$
since the conditions already obtained imply that we do not have any problems of
integrability at $1$.

Now, the inequalities (3.5) are easily  obtained from (3.4) (with $n=1$) and the
estimate (2.1) in a very similar way as it was done in the proof of theorem 2.

In order to verify (1.7), we take the function
$$
f_\lambda (x) = \lambda^{1/2} \chi_{\{W(x)>\lambda^{1/2}\}} (x)
\lambda^{1/2}\chi_{\{|j_0(x)|>\lambda^{1/2}\}} (x)
$$
for each $\lambda >0$,
where $W(x)= \min\{U(x),V(x)^{-1},U(x)V(x)^{-1}\}$. By Holder's inequality and
(3.4),  $f_\lambda$ belongs to $L^p((0,1);V(x)^pxdx) \cap L^2((0,1);xdx)$.
The function $g(x)=f_\lambda (x)\chi_E (x)$, ($E$ measurable $\subset (0,1)$)
satisfies $|g(x)|^p=\lambda^p \chi_{\{|g(x)|>\lambda /2 \}} (x)$ and then, in a
similar way as in the previous proof, we get
$$
\int_E |f_\lambda (x)|^p U(x)^p xdx \leq C\por \int_E |f_\lambda (x)|^p 
V(x)^p xdx, \quad \forall E \subset (0,1),\quad \forall \lambda >0.
$$
Letting $\lambda \rightarrow 0$ and using the dominated convergence
theorem, we obtain
$$
\int_E \chi_{\{W(x)>0\}} (x) U(x)^p x dx \leq C\por
\int_E \chi_{\{W(x)>0\}} (x) V(x)^p x dx,\quad\forall E\subset (0,1)
$$
which implies (1.7).
\medskip
{\hbox{\bf Proof of theorem 4.\enspace}} $(ii)\Rightarrow (i)$. Observing that
$\min \{{1 \over 4},{\alpha +1 \over 2} \}=1/4$ if $\alpha\geq -1/2$,
it is a particular case of theorem 1. In fact, (ii) implies strong boundedness
(that is, boundedness on $L^p$).

 $(i)\Rightarrow (ii)$. By theorem 3, (i) implies (1.5)' and (1.6) and easy
calculations show that the conditions
$$
-1/4\leq {1 \over p}+ {a-1 \over 2} <1/4;\  -1<pb < p-1 \ ; \ -1<pb_k < p-1 \ 
(1\leq k \leq m) \eqno (3.6)
$$
follow from (1.5)' and (1.6).
The strict inequality in ${-1 \over 4}\leq {1 \over p}+ {a-1 \over 2} <{1 \over
4}$ cannot be deduced from conditions (1.5)' and (1.6).

In what follows, we shall use the following decomposition of the kernel
$$
K_n^\alpha(x,y)= K_{n,13}^\alpha(x,y)+K_{n,24}^\alpha(x,y)+K_{n,5}^\alpha(x,y)+
K_{n,6}^\alpha(x,y)
$$
where
$$
K_{n,13}^\alpha(x,y)=K_{n,1}^\alpha(x,y)+K_{n,3}^\alpha(x,y),\quad
K_{n,24}^\alpha(x,y)=K_{n,2}^\alpha(x,y)+K_{n,4}^\alpha(x,y).
$$
So, we can decompose
$$
S_n^\alpha f(x) = S_{n,13}^\alpha f(x) + S_{n,24}^\alpha f(x) +
S_{n,5}^\alpha f(x)+S_{n,6}^\alpha f(x)
$$
where
$$
S_{n,13}^\alpha f(x)=\int_0^1 {M_nJ_\alpha(M_nx)J_{\alpha+1}(M_ny)y \over
y^2 - x^2 } f(y)ydy
$$
and
$$
S_{n,24}^\alpha f(x)=\int_0^1 {M_nJ_\alpha(M_ny)J_{\alpha+1}(M_nx)x \over
x^2 - y^2 } f(y)ydy.
$$

The plan is to prove that under conditions (3.6), $S_{n,24}^\alpha$, 
$S_{n,5}^\alpha$, $S_{n,6}^\alpha$ are bounded from $L^p((0,1);U(x)^pxdx)$ into
$L^p_* ((0,1);U(x)^pxdx)$ (in fact, $S_{n,24}^\alpha$ will be of strong type),
while $S_{n,13}^\alpha$ is not bounded in the case
$-{1 \over 4}={1 \over p}+ {a-1 \over 2}$.

\medskip

\noindent {\it Weak boundedness of $S_{n,5}^\alpha$}. It is simply to observe
that
$$
|S_{n,5}^\alpha f(x)| \leq C |\int_0^1 (xy)^\alpha f(y)ydy|\leq
C x^{-1/2} \int_0^1 y^{-1/2} |f(y)|ydy
$$
and  use Holder's inequality.

\medskip

\noindent {\it Weak boundedness of $S_{n,6}^\alpha$}. It would be enough to
prove (as in the last part of the proof of proposition 1) that
$U(x)^px^{1-p/2}\in A_p(0,1)$, but this is not true, since   the condition 
$-{1 \over 4}<{1 \over p}+ {a-1 \over 2}$ is needed as lemma 4 establishes.

We will proceed as follows: we set the operator $S_{n,6}^\alpha$ in terms of
the operator $J$ and we only need to show 
$$
 \Vert x^{-1/2} \chi_{(0,1/2)}(x)Jf(x)\Vert_{L^p_* ((0,1);U(x)^pxdx)} \leq C
\por \Vert x^{-1/2} f(x)\Vert_{L^p((0,1);U(x)^pxdx)} \eqno (3.7)
$$
$$
\Vert x^{-1/2} \chi_{(1/2,1)}(x)Jf(x)\Vert_{L^p_* ((0,1);U(x)^pxdx)} \leq C
\por \Vert x^{-1/2} f(x)\Vert_{L^p((0,1);U(x)^pxdx)}. \eqno (3.8)
$$
The inequality (3.8)  holds (even with strong norm) because the pair of
weights
$$(U(x)^p x^{1-p/2} \chi_{(1/2,1)}(x),U(x)^p x^{1-p/2})\in
A_p^\delta (0,1)
$$
since only the inequality ${1 \over p}+ {a-1 \over 2} <{1 \over 4}$ is needed in
$x=0$.

For (3.7) notice that if $0<x<1/2$ then $|Jf(x)|\leq C\por \int_0^1 |f(y)|dy$
and
$$\eqalign{&
\Vert x^{-1/2} \chi_{(0,1/2)}(x)Jf(x)\Vert_{L^p_* ((0,1);U(x)^pxdx)} \leq 
\cr
& \leq C \por\int_0^1 |f(y)|dy\por 
\Vert x^{-1/2} \chi_{(0,1/2)}(x)\Vert_{L^p_* ((0,1);U(x)^pxdx)}\leq
\cr
&
\leq C\por \int_0^1 |f(y)|dy\por 
\Vert x^{-1/2}\Vert_{L^p_* ((0,1);U(x)^pxdx)}\leq
\cr
&
\leq C \por \int_0^1 |f(y)|dy\por\leq
 C \Vert x^{-1/2}U(x)f(x)\Vert_{L^p (xdx)}
\Vert x^{-1/2}U(x)^{-1}f(x)\Vert_{L^{p'}(xdx)}\leq
\cr
&
\leq C \Vert x^{-1/2} f(x)\Vert_{L^p((0,1);U(x)^pxdx)},
\cr}$$
where we have used Holder's inequality and that
$$
x^{-1/2}\in L^p_* ((0,1);U(x)^pxdx)\cap L^{p'}((0,1);U(x)^{p'}xdx)
$$
which is immediate from (3.6). The crucial fact is that
$x^{-1/2}\in L^p_*$ even though we have  the equality in 
$-{1 \over 4}\leq {1 \over p}+ {a-1 \over 2}$ (this is not true for strong
$L^p$).

\medskip

\noindent {\it Strong boundedness of $S_{n,24}^\alpha$}. Making the change of
variables $t=y^2,z=x^2$ and taking the function
$g(z)=M_n^{1/2}J_\alpha(M_nz^{1/2})f(z^{1/2})$, the boundedness of 
$S_{n,24}^\alpha$ on $L^p((0,1);U(x)^pxdx)$ turns out to be equivalent to
$$
   \Vert M_n^{1/2}J_{\alpha+1}(M_nz^{1/2})z^{1/2}Hg(z)
   \Vert_{L^p(U(z^{1/2})^pdz)} \leq
$$
$$
   \leq C \Vert M_n^{-1/2}J_\alpha(M_nz^{1/2})^{-1}g(z)
   \Vert_{L^p(U(z^{1/2})^pdz)}.
$$
Now, by using the estimate (3.1), we only need to see that
$$
\Vert Hg(z)\Vert _{L^p(U(z^{1/2})^pz^{p/4}dz)} \leq C\por
\Vert g(z)\Vert _{L^p(U(z^{1/2})^pz^{p/4}dz)}
$$
and this is equivalent to check that $(U(z^{1/2})^pz^{p/4})\in A_p(0,1)$. Since
$1-x^{1/2}\approx 1-x$,\ $|x^{1/2}-x_k|\approx |x-x_k^2|$, the $A_p$-condition
appears as
$$
U(x)=x^{pa/2+p/4}(1-x)^{pb}\prod_{k=1}^m |x-x_k^2|^{pb_k} \in A_p(0,1)
$$ 
which is true from lemma 4 under the hypothesis (3.6).
\medskip
The last step is to prove that, when $-{1 \over 4}={1 \over p}+ {a-1 \over 2}$
there exists no constant $C$ such that
$$
\Vert S_{n,13}^\alpha f(x) \Vert_{L^p_*((0,1);U(x)^pxdx)} \leq
C\por \Vert f(x)U(x) \Vert_p \eqno (3.9)
$$
``uniformly" on $n$.

 The proof is by contradiction. Assume that (3.9) holds, take $0<d<1$ and
$f(x)=\chi_{(d,1)}(x)\por {\rm sgn}(J_{\alpha+1}(M_nx))\por |g(x)|\por x^{-1/2}$,
where $g(x)$ is a suitable function that we will choose below. It is clear that
for $x\in (0,d)$,
$$
\vert S_{n,13}^\alpha f(x) \vert \geq
M_n^{1/2}|J_\alpha(M_nx)|\int_d^1 M_n^{1/2}|J_{\alpha+1}(M_ny)|y^{-1/2}|g(y)|dy.
$$
Then, (3.9) implies
$$
(\int_d^1 M_n^{1/2}|J_{\alpha+1}(M_ny)|y^{-1/2}|g(y)|dy)\por
\Vert M_n^{1/2}J_\alpha(M_nx)\chi_{(0,d)}(x) \Vert_{L^p_*(U(x)^pxdx)}\leq
$$
$$
\leq C\por \Vert \chi_{(d,1)}(x)g(x)x^{-1/2}U(x)\Vert_p
$$
with $C$ independent of $d$, $n$ and $g$. As $\lim\por M_n=+\infty$ we can apply
lemma 5 and its weak version (as it was done in the proof of theorem 3) to the
two factors in the left hand side of the inequality, and we obtain
$$
(\int_d^1 y^{-1}\por |g(y)|dy) \por
\Vert x^{-1/2}\chi_{(0,d)}(x) \Vert_{L^p_*(U(x)^pxdx)} \leq C
\Vert x^{-1/2}\chi_{(d,1)}(x)g(x)U(x) \Vert_p. \eqno (3.10)
$$
Now, let us take $0<d<r<x_1$ (therefore $U(x) \approx C x^a$ on $(0,r)$) and
$g=\chi_{(d,r)}$. Then, it is easy to prove that
$$
\int_d^1 y^{-1} |g(y)|dy=\log {r \over d}, \quad 
\Vert x^{-1/2}\chi_{(d,1)}(x)g(x)U(x) \Vert_p \approx C (\log {r \over d})^{1/p},
$$
and also that if $-{1 \over 4}={1 \over p}+ {a-1 \over 2}$ then
$$
\Vert x^{-1/2}\chi_{(0,d)}(x) \Vert_{L^p_*(U(x)^pxdx)} \approx
C \Vert x^{-1/2}\chi_{(0,d)}(x) \Vert_{L^p_*(x^{ap+1}dx)} \approx C.
$$
Thus, (3.10) implies the existence of a constant $C$ such that
$$
\log {r \over d}\leq C\por  (\log {r \over d})^{1/p}
$$
which is false.

This concludes the proof of theorem 4.

{\hbox{\bf Proof of theorem 5.\enspace}} Let $\alpha\geq -1/2$ and $p=4$.
According to the previous proof, we just need to show the $(4,4)$-restricted
weak type of $S_{n,13}^\alpha$. Making the change of variables $t=y^2,z=x^2$,
this boundedness is equivalent to
$$
|\{z\in (0,1): |M_n^{1/2}J_\alpha(M_nz^{1/2})H(M_n^{1/2}J_{\alpha+1}(M_nt^{1/2})
t^{1/2} f(t^{1/2}))(z)|>\lambda\}| \leq
$$
$$
\leq C\por \lambda^{-4}\int_0^1 |f(x)|^pxdx
$$
where $f$ is a characteristic function. Now, denoting $g(t)=f(t^{1/2})$ and
using that \hfill \break
 $|M_n^{1/2}J_\alpha(M_nz^{1/2})|\leq C\por z^{-1/4}$ (see (3.1)),
it will be enough to prove that
$$
|A_{\lambda ,E}|\leq C\por \lambda^{-4} |E|,\quad \forall \lambda > 0,\quad
\forall E \subseteq  (0,1) \eqno (3.11)
$$
where
$$
A_{\lambda ,E} = \{z\in (0,1):
z^{-1/4}|H(M_n^{1/2}J_{\alpha+1}(M_nt^{1/2}) t^{1/2}\chi_E(t))(z)|>\lambda\}.
$$
Observe that
$$
|A_{\lambda ,E}\cap (1/4,1)|\leq
C\por \lambda^{-4}\int_0^1
|H(M_n^{1/2}J_{\alpha+1}(M_nt^{1/2}) t^{1/2}\chi_E(t))(z)|^4dz\leq
$$
$$
\leq C\por \lambda^{-4}\int_0^1
|M_n^{1/2}J_{\alpha+1}(M_nt^{1/2}) t^{1/2}\chi_E(t)|^4dt \leq
$$
$$
\leq C\por \lambda^{-4}\int_0^1 t^{1/4}|\chi_E(t)|^4
\leq C\por \lambda^{-4}|E|
$$
due to the boundedness of the Hilbert transform on $L^4((0,1);dx)$ and to the
estimate (3.1).

On the other hand, if $E\subseteq [1/2,1)$ and $z\in (0,1/4)$
$$
|H(M_n^{1/2}J_{\alpha+1}(M_nt^{1/2}) t^{1/2}\chi_E(t))(z)|\leq C
\int_0^1 t^{1/4}\chi_E(t)dt \leq C \por |E|.
$$
Hence,
$$
|A_{\lambda ,E}\cap (0,1/4)|\leq |\{z\in (0,1/4): z^{-1}\por C \por |E|^4> 
\lambda^4\}|\leq C\por \lambda^{-4}|E|^4 \leq C\por \lambda^{-4}|E|
$$
and (3.11) holds, if $E\subseteq [1/2,1)$.

If $E\subseteq (0,1/2)$, let $I_k=({1\over 2^{k+1}},{1\over 2^{k}})$,
$k=2,3,...$. We shall use the following notation:
$$
I_{k1}=(0,{1\over 2^{k+2}}),\quad  I_{k2}=({1\over 2^{k+2}},{1\over 2^{k-1}}),
\quad  I_{k3}=({1\over 2^{k-1}},1/2),\quad k=2,3,...
$$
$$
A_{\lambda ,E}^k=A_{\lambda ,E}\cap I_k,\quad
A_{\lambda ,E}^{km}=A_{\lambda /3 ,E_{km}}^k,\ m=1,2,3
$$
where $E_{km}=E\cap I_{km}$, $m=1,2,3$.

Then, $A_{\lambda ,E}^k\subseteq A_{\lambda ,E}^{k1}\cup A_{\lambda ,E}^{k2}\cup
A_{\lambda ,E}^{k3}$ and
$$
|A_{\lambda ,E}\cap (0,1/4)|\leq \sum_{k=2}^\infty |A_{\lambda ,E}^k |\leq
\sum_{k=2}^\infty \{|A_{\lambda ,E}^{k1}|+|A_{\lambda ,E}^{k2}|+|A_{\lambda
,E}^{k3}|\}. \eqno (3.12)
$$
For $m=3$, notice that if $z\in I_k$ and $t\in I_{k3}$ then $t-z\approx t$.
Hence, using (3.1)
$$
|H(M_n^{1/2}J_{\alpha+1}(M_nt^{1/2}) t^{1/2}\chi_{E_{k3}}(t))(z)|\leq 
C\por\int_0^1 t^{-3/4}\chi_E(t)dt \leq C\por |E|^{1/4}
$$
(the last inequality can be deduced from duality between Lorentz spaces).
Therefore,
$$
|A_{\lambda ,E}^{k3}|\leq |\{z\in I_k : z< C\por \lambda^{-4}|E|\}|. \eqno
(3.13)
$$
The same arguments work in order to prove that  $|A_{\lambda ,E}^{k1}|$ has the
estimate (3.13).

For $m=2$, we have
$$\eqalign{
|A_{\lambda ,E}^{k2}|&\leq |\{z\in I_k :
|H(M_n^{1/2}J_{\alpha+1}(M_nt^{1/2})
t^{1/2}\chi_{E_{k2}}(t))(z)|>C\por\lambda\por 2^{-{k\over 4}}\}|\cr
&\leq C\por \lambda^{-4}\por 2^{k}
\int_0^1 |H(M_n^{1/2}J_{\alpha+1}(M_nt^{1/2}) t^{1/2}\chi_{E_{k2}}(t))(z)|^4dz
\cr & \leq C\por \lambda^{-4}\por 2^{k}
\int_0^1 |M_n^{1/2}J_{\alpha+1}(M_nt^{1/2}) t^{1/2}\chi_{E_{k2}}(t)|^4dt
\cr & \leq C\por \lambda^{-4}\por 2^{k} \int_0^1 t \chi_{E_{k2}}(t)
\leq C\por \lambda^{-4} |E_{k2}|.
\cr}\eqno (3.14)$$

Finally, putting together (3.12), (3.13) and (3.14) we get (3.11) and the proof
for $p=4$ is concluded.

The proof for $p=4/3$ follows from standard arguments of duality between
Lorentz spaces.
\medskip
\noindent {\bf Remarks.} The decomposition in four terms which we have used to
study the weak boundedness,  also works to study the strong boundedness,
instead of using the decomposition in six terms. However, the results are
somehow different. For instance, to obtain the uniform boundedness of the
proposition 1, we must change the condition by 
$$
(U(x^{1/2})^px^{-p/4},V(x^{1/2})^px^{-p/4})\in A_p^\delta(0,1),\ \ 
(U(x^{1/2})^px^{p/4},V(x^{1/2})^px^{p/4})\in A_p^\delta(0,1).
$$
Something similar happens with proposition 2. If we apply these conditions to
particular weights as in theorem 1, the inequalities (1.3), (1.4) are the same.

Analogous results can be obtained in a similar way for Fourier-Dini series.
\bigskip
\noindent {\bf References.}\par\nobreak
\medskip
\item {{\bf 1.}} A. I. Benedek and R. Panzone, Mean convergence of series of
Bessel functions, {\it Rev. Un. Mat. Arg.} {\bf 26} (1972), 42--61.

\item {{\bf 2.}} S. Chanillo, On the weak behaviour of partial sums of
Legendre series, {\it Trans. Amer. Math. Soc.} {\bf 268} (1981), 367--376.

\item {{\bf 3.}} S. Chanillo, The multiplier for the ball and radial
functions, {\it J. Funct. Anal.} {\bf 55} (1984), 18--24.

\item {{\bf 4.}} J. Garc\'\i a Cuerva and J. L. Rubio de Francia, \lq \lq
Weighted norm inequalities and related topics", North-Holland, Amsterdam, 1985.

\item {{\bf 5.}} J. J. Guadalupe, M. P\'erez, F. J. Ruiz and J. L. Varona,
Two notes on convergence and divergence a.e. of Fourier series with
respect to some orthogonal systems, (preprint).

\item {{\bf 6.}} J. J. Guadalupe, M. P\'erez and J. L. Varona, Weak behaviour of
Fourier-Jacobi series, {\it J. Approx. Theory} {\bf 61} (1990), 222--238.

\item {{\bf 7.}} R. Hunt, B. Muckenhoupt and R. Wheeden, Weighted norm
inequalities for the conjugate function and Hilbert transform, {\it Trans. Amer.
Math. Soc.} {\bf 176} (1973), 227--251.

\item {{\bf 8.}} A. M\'at\'e, P. Nevai and V. Totik, Necessary
conditions for weighted mean convergence of Fourier series in orthogonal
polynomials, {\it J. Approx. Theory} {\bf 46} (1986), 314--322.

\item {{\bf 9.}} B. Muckenhoupt, Weighted norm inequalities for the
Hardy maximal function, {\it Trans. Amer. Math. Soc.} {\bf 165} (1972), 207--226.

\item {{\bf 10.}} C. J. Neugebauer, Inserting $A_p$ weights, {\it Proc. Amer.
Math. Soc.} {\bf 87} (1983), 644--648.

\item {{\bf 11.}} E. Romera and F. Soria, Endpoint estimates for the
maximal operator associated to spherical partial sums on radial functions,
(preprint).

\item {{\bf 12.}} G. N. Watson, \lq \lq A treatise on the theory of
Bessel functions", 2nd. ed., Cambridge Univ. Press, Cambridge, 1958.

\item {{\bf 13.}} G. M. Wing, The mean convergence of orthogonal series,
{\it Amer. J. Math.} {\bf 72} (1950), 792--808.

\bigskip
\noindent {\rmdiez Jos\'e J. Guadalupe, Mario P\'erez and Francisco
J. Ruiz. Departamento de Matem\'aticas. Facultad de Ciencias. Universidad de
Zaragoza, 50009 Zaragoza, SPAIN.}

\noindent {\rmdiez Juan L. Varona. Departamento de Matem\'atica
Aplicada. Colegio Universitario de la Rioja, 26001 Logro\~no, SPAIN.}

\bye
